Two observations follow. First, the result is not explained by model
size alone: Qwen2-VL-2B and LLaVA-Interleave-0.5B received the same
input, and only the larger of the two used it. Model size cannot be
compared with PaliGemma or BLIP-2 because their inputs were shorter.
Second, the full-card condition is not a pure test of visual or spectral
understanding. Qwen2-VL-2B may be reading the summary line, the index
values, the images or a mixture of these, and it has not yet been tested
on the full test set without the summary line. For LLaVA, removing the
line made almost no difference (32.7\% to 36.4\%), because it did not
use the card in either condition.

The baselines show that the card contains enough information to solve
the task: a single threshold on its $\Delta$NDVI value reaches 94.5\%
with cut-offs fitted on the training set. This is expected, since the
labels are defined from NDVI change, and it sets a clear target that
none of the VLMs approaches. The closed-set probe lifts LLaVA to 60.0\%,
which suggests that constraining the output helps a small model even
when free generation fails, although the probe never predicts
``stable''.

\subsection{Effect of Fine-Tuning}

The class LoRA raises Qwen2-VL-2B from 56.4\% to \textbf{67.3\%} (37/55;
95\% interval 54.1--78.2\%; Table~\ref{tab:lora}). The intervals of the
two runs overlap, so with 16 training patches and $n = 55$ test patches
this is an indicative gain rather than a powered comparison. The
fine-tuned model was correct on all 18 decline and all 19 increase test
patches, but it labelled every stable patch as decline. The adapter
therefore learned two of the three classes and lost the third.

\begin{table}[t]
\centering
\small
\begin{tabular}{@{}lcc@{}}
\toprule
Qwen2-VL-2B run & Train & Test acc. \\
\midrule
Zero-shot (full card) & 0 & 56.4\% \\
Class LoRA, 1 epoch (full card) & 16 & 67.3\% \\
\midrule
\multicolumn{3}{l}{\emph{Excluded or exploratory}} \\
Full card, 2--3 epochs & 256 & 100\% \\
Card in training, images-only test & 128 & 32.7\% \\
\bottomrule
\end{tabular}
\caption{Class LoRA results for Qwen2-VL-2B. The 100\% runs copied the
label from the summary line and are not counted as results.}
\label{tab:lora}
\end{table}

Two exploratory runs support this reading of the summary line. Adapters
trained on all 256 patches with the full card reached 100\% because
they could read the answer from the card. An adapter trained with the
card on 128 patches and then tested on images alone scored 32.7\%,
below the majority baseline. Neither run is a valid estimate of
performance, and the second mixes training and test conditions, but
together they suggest that the fine-tuned model relies on the card text
rather than on the images. The explanation LoRA is a
different task and is discussed in Section~\ref{sec:expl}.

\subsection{Error Analysis}

Zero-shot Qwen2-VL-2B was correct on all 18 decline patches and on 13 of
18 stable patches, but it answered ``stable'' for all 19 increase
patches. For about half of the increase class the card's trend words
say ``stable'', and the model followed those words rather than the
summary line. Errors therefore concentrate where the label is a weak
rank: increase patches with $\Delta$ near zero are reasonably called
stable, and decline and stable patches overlap slightly after the
second-batch percentile shift.

The 108 cards whose trend words disagree with the label build a
contradiction into the input: a model that reads the descriptive text
faithfully can still miss the gold label. Together with the low human
agreement (8 of 19), this means that some ``errors'' are likely to be
label noise rather than model failure. The test set was not relabelled
after seeing model outputs.

The failures of the other models are mostly failures to follow the task.
LLaVA always produced the same answer, PaliGemma refused long prompts
and largely answered with one class, and BLIP-2 echoed its instruction.
The parser accepted the echo because it searched the whole output for a
class word; a stricter parser should read only newly generated text and
reject copied instructions.

\subsection{Explanation Quality}
\label{sec:expl}

All explanation runs used the card without the label line.
Table~\ref{tab:explain} summarises them. With the first, free-form
prompt, which asked the model to describe and justify the change in its
own words, PaliGemma answered with a class or ``unanswerable'' and no
reason, and BLIP-2 repeated the prompt, so neither produced a checkable
claim. LLaVA wrote fluent but repetitive text in which 56.7\% of the
checkable claims were correct. Qwen2-VL-2B was run on only three
patches in this mode, and one of its claims misstated the rainfall
change. Free-form output from these models is not usable as an audit
trail.

\begin{table*}[t]
\centering
\small
\begin{tabular}{lrrrrr}
\toprule
Setting and model & Class acc. & Claims / expl. & Correct claims & $\ge 1$ wrong & Self-consistent \\
\midrule
Free-form, LLaVA-Interleave-0.5B & 29.1\% & 1.09 & 56.7\% & --- & --- \\
Free-form, PaliGemma-3B & 33.3\%$^{a}$ & 0 & --- & --- & --- \\
Free-form, BLIP-2 OPT-2.7B & 32.7\%$^{b}$ & 0 & --- & --- & --- \\
\midrule
Structured, Qwen2-VL-2B & 63.6\% & 2.85 & \textbf{75.8\%} & 36.4\% & 100\% \\
Structured, LLaVA-Interleave-0.5B & 45.1\%$^{c}$ & 2.33 & \textbf{71.1\%} & 34.5\% & 90.7\% \\
\midrule
Explanation LoRA, Qwen2-VL-2B & 61.8\% & 6.84 & 50.5\% & 74.5\% & 96.4\% \\
\bottomrule
\end{tabular}
\caption{Explanation results on the test set, label removed from the
card. $^{a}$12 of 36 answered patches. $^{b}$Prompt echo; every answer
parsed as ``stable''. $^{c}$23 of 51 answered patches. The free-form
Qwen2-VL-2B run (3 patches, 43.8\% correct claims) is omitted.}
\label{tab:explain}
\end{table*}

The structured prompt gives the model a numbered list of allowed facts
that must be copied verbatim as supporting evidence, a line stating an
NDVI decision rule for the patch ($\pm 0.02$), and a request for a
one-sentence justification about NDVI only. Under this prompt, claim
correctness rises to 75.8\% for Qwen2-VL-2B and 71.1\% for
LLaVA-Interleave-0.5B. Most of the remaining incorrect claims are not
invented facts: 35 of Qwen's 38 and 35 of LLaVA's 37 flagged claims
describe a copied change smaller than the scorer's dead-band (for
example, an EVI change of $-0.006$ described as a decrease). High claim
correctness here means that the model did not invent numbers; it does
not mean that the model reasoned about the change
. Qwen2-VL-2B's class matched the stated
rule on 53 of 55 patches, and the rule alone scores 63.6\%, exactly
Qwen's accuracy in this mode. This is instruction-following, not a new
classification ability.

The explanation LoRA reached 61.8\% class accuracy (34/55; 95\% interval
48.6--73.5\%) and produced more claims per explanation (6.84), but only
50.5\% of them were correct and 74.5\% of explanations contained at
least one incorrect claim. This is worse than structured zero-shot
prompting. The bottleneck is the teacher: only 28 of 256 of its own
generations passed the filter, and only 16 were used. With this
evidence, self-distillation of a 2B model does not teach faithful
explanation; the adapter mainly learned the output format.

\subsection{Comparison with Existing Approaches}
